% Hitoshi Inamori
% Cheat sheets Standard Format File


%\documentstyle[11pt]{article}
\documentclass[11pt,amsmath,amssymb,latexsym]{article}



% Default margins are too wide all the way around.  I reset them here
\setlength{\topmargin}{-.5in}
\setlength{\textheight}{9in}
\setlength{\oddsidemargin}{-0.3in}
%\setlength{\evensidemargin}{.125in}
\setlength{\textwidth}{6.25in}


\usepackage{leftidx}% http://ctan.org/pkg/leftidx
\usepackage{slashed}

%%%%%%%%%%%%%%%%%%%%%%%%%%%%%%%%%%%%%%%%%%%%%%%%%%%%%%%%%%%%%%%%%%%%%%%%%%%
% My commands are here: basically we create a new standard table using
% \toshTable
% Then 
%		a definition is created using \toshDefinition
%   a property is created using \toshProperty
%   a proof (one column) is created using \toshProof
%   an example (one column) is created using \toshExample
%%%%%%%%%%%%%%%%%%%%%%%%%%%%%%%%%%%%%%%%%%%%%%%%%%%%%%%%%%%%%%%%%%%%%%%%%%%
\newcommand{\toshTableBegin} 			{\begin{tabular}{|p{5cm}|p{11cm}|}}
\newcommand{\toshTableEnd} 				{\hline \end{tabular}}
\newcommand{\toshTableTitle}[1] 	{\hline \multicolumn{2}{|l|}{{\bf{#1}}}\\ \hline}
\newcommand{\toshTableNotation}[1]{\hline \multicolumn{2}{|p{16cm}|}{(\textbf{Note}) #1}\\}
\newcommand{\toshDefinition}[2]		{\hline (D) {\bf{#1}} & {#2} \\}
\newcommand{\toshProperty}[2]			{\hline (P) {#1} & {#2} \\}
\newcommand{\toshProof}[1]				{\hline \multicolumn{2}{|p{16cm}|}{(\small \underline{Proof}) #1}\\}
\newcommand{\toshExample}[1]			{\hline \multicolumn{2}{|p{16cm}|}{\small (Example) #1}\\}
\newcommand{\toshNewPage}					{\hline \end{tabular} \newpage \toshTableBegin}
%%%%%%%%%%%%%%%%%%%%%%%%%%%%%%%%%%%%%%%%%%%%%%%%%%%%%%%%%%%%%%%%%%%%%%%
% For Definitions or Properties with several conditions or implications
%%%%%%%%%%%%%%%%%%%%%%%%%%%%%%%%%%%%%%%%%%%%%%%%%%%%%%%%%%%%%%%%%%%%%%%
\newcommand{\toshBlockStart}		{\begin{tabular}{l}}
\newcommand{\toshBlockItem}[1]	{{#1} \\}
\newcommand{\toshBlockDesc}[2]	{{\bf{(#1)}}\, {#2} \\}
\newcommand{\toshBlockEnd}			{\end{tabular}}
\newcommand{\toshLeftBlock}[1]	{$\left.{#1}\right]$}
\newcommand{\toshRightBlock}[1]	{$\left[{#1}\right.$}
\newcommand{\toshDef}[1]				{{\textbf{#1}}}
\newcommand{\toshDual}{\tilde{E}}
%%%%%%%%%%%%%%%%%%%%%%%%%%%%%%%%%%%%%%%%%%%%%%%%%%%%%%%%%%%%%%%%%%%%%%%
% Preferred Symbols
%%%%%%%%%%%%%%%%%%%%%%%%%%%%%%%%%%%%%%%%%%%%%%%%%%%%%%%%%%%%%%%%%%%%%%%
\newcommand{\toshEquivalent}		{$\iff$}
\newcommand{\toshImply}					{\Rightarrow}
\newcommand{\toshByDef}					{$\stackrel{D}{\iff}$}
\newcommand{\toshEqByDef}				{\,\stackrel{D}{=}\,}
\newcommand{\toshIdentity}			{{\mathbf{1}}}
\newcommand{\toshItoInt}				{{\mathcal{I}}}
\newcommand{\toshB}							{{\mathcal{B}}\,}
\newcommand{\toshC}							{{\mathbf{C}}}
\newcommand{\toshD}							{{\mathcal{D}}}
\newcommand{\toshCylinder}			{{\mathcal{C}}_1}
\newcommand{\toshF}							{{\mathcal{F}}\,}
\newcommand{\toshG}							{{\mathcal{G}}\,}
\newcommand{\toshH}							{{\mathcal{H}}\,}
\newcommand{\toshL}							{{\mathcal{L}}}
\newcommand{\toshLOne}					{{\mathcal{L}}^1}
\newcommand{\toshM}							{{\mathcal{M}}\,}
\newcommand{\toshMStep}					{M^2_{{step}}}
\newcommand{\toshMTwo}					{M^2}					
\newcommand{\toshN}							{{\mathbf{N}}\,}
\newcommand{\toshNuST}					{{\nu^{m\times n}_{\toshH}(S,T)}}
\newcommand{\toshST}						{\,/\,} % Such That
\newcommand{\toshR}							{{\mathbf{R}}}
\newcommand{\toshT}							{{\mathbf{T}}}
\newcommand{\toshVar}						{{\textrm{Var}}}

\newcommand{\grad}							{{\mathbf{grad}\,}}
\newcommand{\divergence}						{{\textrm{div}\,}}
\newcommand{\rot}							{{\mathbf{rot}\,}}
%%%%%%%%%%%%%%%%%%%%%%%%%%%%%%%%%%%%%%%%%%%%%%%%%%%%%%%%%%%%%%%%%%%%%%%
% Special Relativity and Electromagnetism
%%%%%%%%%%%%%%%%%%%%%%%%%%%%%%%%%%%%%%%%%%%%%%%%%%%%%%%%%%%%%%%%%%%%%%%
\newcommand{\state}[1]{\left|{#1}\right>}
\newcommand{\stateBasis}[2]{\left|{#1}\right>_{#2}}
\newcommand{\bra}[1]{\left<{#1}\right|}
\newcommand{\ket}[1]{\left|{#1}\right>}
\newcommand{\braBasis}[2]{\left<{#1}\right|_{#2}}
\newcommand{\braket}[2]{\left<{#1}|{#2}\right>}
%\newcommand{\braketBasis}[3]{\leftidx{_{#3}}{\left<{#1}|{#2}\right>}_{#3}}
\newcommand{\Tr}[1]{\textrm{Tr}\left({#1}\right)}
\newcommand{\PartTr}[2]{\textrm{Tr}_{#2}\left({#1}\right)}


%%%%%%%%%%%%%%%%%%%%%%%%%%%%%%%%%%%%%%%%%%%%%%%%%%%%%%%%%%%%%%%%%%%%%%%
% Begin document
%%%%%%%%%%%%%%%%%%%%%%%%%%%%%%%%%%%%%%%%%%%%%%%%%%%%%%%%%%%%%%%%%%%%%%%
\begin{document}
\title{Quantum field theory cheat sheet\\ -- path-integral formulation -- }
\author{Hitoshi Inamori}
%\renewcommand{\today}{April, 2008}
\maketitle
\begin{abstract}
From Peskin Schroeder ``An introduction to Quantum Field Theory Chapter 9''
\end{abstract}


%%%%%%%%%%%%%%%%%%%%%%%%%%%%%%%%%%%%%%%%%%%%%%%%%%%%%%%%%%
% Green function
%%%%%%%%%%%%%%%%%%%%%%%%%%%%%%%%%%%%%%%%%%%%%%%%%%%%%%%%%%

\toshTableBegin
\toshTableTitle{Green function in path integral formulation}
\toshProperty{$n$-point Green function or correlation function}{The probability amplitude to observe $k_1, k_2, \ldots, k_n$ particles described by a field $\phi$ at space-time positions $x_1, x_2, \ldots, x_n$ is given by the \toshDef{$n$-point Green function} or \toshDef{$n$-point correlation function} \underline{denoted} as $\bra{\Omega} T\left[ \phi_H^{k_1}(x_1)\phi_H^{k_2}(x_2)\cdots\phi_H^{k_n}(x_n)\right]\ket{\Omega}$. 

\smallskip

According to the path-integral formulation, this function can be obtained from a Lorentz invariant Lagrangian $\toshL$:

\bigskip

\quad \framebox{\quad$\displaystyle \frac{\displaystyle\int\toshD\phi\, \phi^{k_1}(x_1)\phi^{k_2}(x_2)\cdots\phi^{k_n}(x_n) \exp\left[i\int d^4 x \toshL(\phi,\partial_\mu \phi)\right]}{\displaystyle\int\toshD\phi\, \exp\left[i\int d^4 x \toshL(\phi,\partial_\mu \phi)\right]}$\quad}

\bigskip

Note that in the above formula, there are no operators involved.

\smallskip

Note that each degree of freedom of $\phi$ (as a real number, i.e. a complex field has 2 degrees of freedom) represents a degree of freedom for the particle (spin, particle/anti-particle). 

\smallskip

Note that the observed fields must be explicitly listed. For instance, absence of $\phi^{k_2}(x_2)$ above means that there is no particle at point $x_2$, which is not equivalent to saying that we do not observe what happens at $x_2$.
}
\toshTableEnd
\bigskip

%%%%%%%%%%%%%%%%%%%%%%%%%%%%%%%%%%%%%%%%%%%%%%%%%%%%%%%%%%
% Functional derivatives, generating function
%%%%%%%%%%%%%%%%%%%%%%%%%%%%%%%%%%%%%%%%%%%%%%%%%%%%%%%%%%

\toshTableBegin
\toshTableTitle{Functional derivatives, generating function}
\toshDefinition{Functional derivatives for fields}{
If space-time was discretized and each space-time point indexed by $i$, a field would be a vector $(\phi_i)_i$ and if we have a functional on the field $\phi$, a derivative with respect to $\phi_i$ for some space-time point indexed by $i$ would be defined as a standard derivative with respect to $\phi_i$. Namely $\frac{\partial}{\partial x_i} x_j = \delta_{i j}$ and $\frac{\partial}{\partial x_i} \sum_j x_j k_j = k_j$.

\smallskip

The \toshDef{functional derivative} is an extension of this for fields defined on a continous space-time:

\quad $\displaystyle \frac{\delta}{\delta J(x)} J(y) = \delta(x-y)$, \quad and \quad $\displaystyle \frac{\delta}{\delta J(x)} \int d^4 y J(y)\psi(y) = \psi(x)$

\smallskip

For composite functions, the ordinary differentiation rules applies:

\quad $\displaystyle \frac{\delta}{\delta J(x)} f(G[ J]) =  \frac{\delta G[J] }{\delta J(x)} f'[G[J]]$, for instance

\quad $\displaystyle \frac{\delta}{\delta J(x)} \exp\left[i\int d^4 y J(y)\psi(y)\right] = i\psi(x)\exp\left[i\int d^4 y J(y)\psi(y)\right]$

\smallskip

For functional on derivatives of $J$, first integrate by parts to remove the derivatives (assuming surface terms are zero), for instance:

\quad $\displaystyle \frac{\delta}{\delta J(x)} \int d^4 y (\partial_\mu J(y)) V^\mu(y) = - \partial_\mu V^\mu$

}
\toshDefinition{Generating functional}{The correlation functions can be written as functional derivatives of a functional $Z[J]$ called \toshDef{generating functional}:

\quad $\displaystyle Z[J] = \int\toshD\phi \exp\left[ i\int d^4 x [\toshL(x) + J(x)\phi(x)]\right]$

for instance

 $\displaystyle \bra{\Omega} T\left[ \phi_H(x_1)\phi_H(x_2) \right]\ket{\Omega} = \frac{1}{Z_0} \left.\left(-i\frac{\delta}{\delta J(x_1)}\right)\left(-i\frac{\delta}{\delta J(x_2)}\right) Z[J]\right|_{J=0}$

where $Z_0 = Z[J=0]$.
}
\toshTableEnd
\bigskip

%%%%%%%%%%%%%%%%%%%%%%%%%%%%%%%%%%%%%%%%%%%%%%%%%%%%%%%%%%
% Free real scalar fields
%%%%%%%%%%%%%%%%%%%%%%%%%%%%%%%%%%%%%%%%%%%%%%%%%%%%%%%%%%

\toshTableBegin
\toshTableTitle{Free real scalar fields}
\toshProperty{Generating functional}{The free Lagrangian of the real scalar field is 

\quad $\displaystyle \toshL_0 = \frac{1}{2}(\partial_\mu\phi)(\partial^\mu\phi) - \frac{1}{2}m^2\phi^2$

The exponent for the generating functional reads, after an integration by part and adding a $i\epsilon$ component so that the integral converges:

\quad $\displaystyle \int d^4 x [\toshL_0 + J\phi] = \int d^4 x\left[\frac{1}{2}\phi(-\partial^2- m^2 + i\epsilon)\phi + J\phi\right]$

using the change of variable $\phi(x)=\phi'(x)+ i\int d^4 y \Delta_F(x-y)J(y)$ with Jacobian 1 (where $\Delta_F$ is the Feynman propagator), using a double integration by parts and using the fact that $\Delta_F(x-y)$ is a Green function for the operator $(\partial^2+m^2)$ we get the generating functional:

\smallskip

\quad $\displaystyle Z[J] = Z_0 \exp\left[ -\frac{1}{2} \int d^4 x d^4 y \, J(x) \Delta_F(x-y)J(y)\right]$

\smallskip

we check in particular that

\quad $\displaystyle \bra{0} T[\phi(x_1)\phi(x_2)]\ket{0} = -\frac{\delta^2}{\delta J(x_1) \delta J(x_2)} Z[J] = \Delta_F(x_1-x_2)$

}
\toshTableEnd
\bigskip


%%%%%%%%%%%%%%%%%%%%%%%%%%%%%%%%%%%%%%%%%%%%%%%%%%%%%%%%%%
% Free electromagnetic fields
%%%%%%%%%%%%%%%%%%%%%%%%%%%%%%%%%%%%%%%%%%%%%%%%%%%%%%%%%%

\toshTableBegin
\toshTableTitle{Free electromagnetic fields}
\toshProperty{Generating functional}{The free Lagrangian for the electromagnetic field is

\quad $\toshL_0 = -\frac{1}{4}F^{\mu\nu}F_{\mu\nu}$, where $F_{\mu\nu} = \partial_\mu A_\nu - \partial_\nu A_\mu$

To deal with the gauge invariance, i.e. invariance under the shift 

\quad $A_\mu(x) \mapsto A_\mu(x) + \frac{1}{e}\partial_\mu \alpha(x)$

we restrict the functional integration to the field $A_\mu$ that nearly obeys the Lorentz gauge $G_\omega(A) = \partial^\mu A_\mu (x) - \omega(x) = 0$ with $\omega$ small. In practice we integrate over the functions $\alpha$ such that $G(A)=0$ and integrate over the function $\omega$ with the weight $\exp\left[-i\int d^4 \frac{\omega^2}{2\xi}\right]$. This gives, noting that the action $S[A]$ and any observable $O(A)$ being gauge invariant:

\quad $\displaystyle\int \toshD A e^{iS[A]}$

\quad $\displaystyle = \int\toshD\omega e^{-i\int d^4 x \frac{\omega^2}{2\xi}} \det(G_\omega)\int\toshD\alpha \int\toshD A e^{iS[A]} \delta(G_\omega(A))$

\quad $= \det(G_\omega)\int\toshD\alpha \int\toshD A e^{iS[A]} \exp\left[-i\int d^4 x\frac{1}{2\xi}(\partial^\mu A_\mu)^2\right]$

where $\det(G_\omega) = \det(\partial^2/e)$ is independent of $\alpha$ and $A$. We get:

\quad $\displaystyle\bra{\Omega} T[O(A)]\ket{\Omega} = \frac{\displaystyle \toshD A O(A) \exp\left[ i\int d^4 x [\toshL_0 - \frac{1}{2\xi}(\partial^\mu A_\mu)^2]\right]}{\displaystyle \toshD A \exp\left[ i\int d^4 x [\toshL_0 - \frac{1}{2\xi}(\partial^\mu A_\mu)^2]\right]}$

}
\toshTableEnd
\bigskip

%%%%%%%%%%%%%%%%%%%%%%%%%%%%%%%%%%%%%%%%%%%%%%%%%%%%%%%%%%
% Grassmann numbers
%%%%%%%%%%%%%%%%%%%%%%%%%%%%%%%%%%%%%%%%%%%%%%%%%%%%%%%%%%

\toshTableBegin
\toshTableTitle{Grassmann numbers}
\toshDefinition{Grassmann numbers}{By definition, if $\theta$ and $\eta$ are \toshDef{Grassmann numbers} and $a$ an ordinary number

\smallskip

\toshBlockStart
\toshBlockDesc{anticommutation} {$\theta \eta = - \eta\theta$ thus $\theta^2=0$}
\toshBlockDesc{integral, real number} {$\int d\theta \, a = 0$, $\int d\theta \, \theta = 1$}
\toshBlockDesc{integral, complex number} {$\theta$ and $\theta^*$ are considered as}
\toshBlockItem{\quad  independent numbers, and $\int d\theta^* d\theta (\theta\theta^*) = 1$}
\toshBlockDesc{derivative} {$\frac{\partial}{\partial \eta} \theta\eta = - \frac{\partial}{\partial \eta} \eta\theta = -\theta$}
\toshBlockEnd

\smallskip

Addition, multiplication with ordinary numbers follow ordinary rules, in particular $a\theta = \theta a$.
}
\toshProperty{Properties}{
For any Grassman numbers $\theta$, and ordinary numbers $b$, and Hermitian matrix $B$:

\toshBlockStart
\toshBlockDesc{1} {any function of $\theta$ reads $f(\theta) = a + b\theta$ (Taylor expansion)}
\toshBlockDesc{2} {in an integration $\int d\theta_1 d\theta_2\cdots$ only the integrand $\theta_1\theta_2\cdots$}
\toshBlockItem{\quad contributes}
\toshBlockDesc{3} {$\int d\theta^* d\theta\, e^{-\theta^* b \theta} = b$}
\toshBlockDesc{4} {$\int d\theta^* d\theta\, \theta\theta^* e^{-\theta^* b \theta} = 1$}
\toshBlockDesc{5} {$\left(\prod_i \int d\theta^*_i d\theta_i\right) e^{-\theta_i B_{i j} \theta_j} = \det B$ }
\toshBlockDesc{6} {$\left(\prod_i \int d\theta^*_i d\theta_i\right) \theta_k \theta_l e^{-\theta_i B_{i j} \theta_j} = \det B (D^{-1})_{kl}$ }
\toshBlockEnd

}
\toshTableEnd
\bigskip

%%%%%%%%%%%%%%%%%%%%%%%%%%%%%%%%%%%%%%%%%%%%%%%%%%%%%%%%%%
% Free spinor fields
%%%%%%%%%%%%%%%%%%%%%%%%%%%%%%%%%%%%%%%%%%%%%%%%%%%%%%%%%%

\toshTableBegin
\toshTableTitle{Free spinor fields}
\toshProperty{Generating functional}{The free Lagrangian for the spinor field is

\quad $\displaystyle \toshL_0 = \overline{\psi} (i\slashed{\partial} - m)\psi$

The spinor field is a Grassmann field i.e. it reads as

\quad $\psi(x) = \sum_i \psi_i \phi_i(x)$

where $\phi_i(x)$ is the basis of four-component spinors made of ordinary complex numbers and $\psi_i$ are constant Grassmann complex numbers.

\smallskip

The generating functional reads

\quad $\displaystyle Z[\overline{\eta},\eta ] = \int\toshD\overline{\psi}\toshD\psi \, \exp\left[ i \int d^4 x \left[\overline{\psi} (i\slashed{\partial}-m)\psi + \overline{\eta}\psi + \overline{\psi}\eta\right]\right]$

\smallskip

\quad \quad \quad $\displaystyle = Z_0 \exp\left[ - \int d^4 x \, d^4 y \, \overline{\eta}(x)S_F(x-y)\eta(y)\right]$


}


\toshTableEnd
\bigskip

%%%%%%%%%%%%%%%%%%%%%%%%%%%%%%%%%%%%%%%%%%%%%%%%%%%%%%%%%%
% Quantum electrodynamics
%%%%%%%%%%%%%%%%%%%%%%%%%%%%%%%%%%%%%%%%%%%%%%%%%%%%%%%%%%

\toshTableBegin
\toshTableTitle{Quantum electrodynamics}
\toshProperty{Generating functional}{The full Lagrangian reads

\quad $\displaystyle \toshL = \overline{\psi}(i\slashed{\partial} - m)\psi - \frac{1}{4} F_{\mu\nu} F^{\mu \nu} - e\overline{\psi}\gamma^\mu A_\mu\psi$

\quad \quad $\displaystyle = \toshL_0 - e\overline{\psi}\gamma^\mu A_\mu\psi$

The generating function for the full theory can be written perturbatively expanding the exponential for the interaction term:

\smallskip

$\displaystyle\exp\left[ i \int d^4 x \toshL \right] = \exp\left[ i \int d^4 x \toshL_0 \right] \times\left\{ 1 - ie\int d^4 x \overline{\psi}\gamma^\mu A_\mu \psi + \cdots\right\}$

}


\toshTableEnd
\bigskip

 
\end{document}
